\begin{annexitem}{annex:p12.ue.forcing}{thm:p12-ue-forcing}{The elementary binding calculation on the entire alphabet}
\begin{proof}
First, simultaneous composition acts on terms as
\(s(\theta\eta)=(s\theta)\eta\). The variable clause is precisely
the composition definition, constants are unchanged, and a function
application follows by induction on its finite argument list. Applying
the term law to variable images proves associativity; it also proves
finite support of composition, whose support is contained in the union
of the two supports. No successive rewriting of an inserted image is used.

If \(\rho\) solves \(x\doteq t\), then
\(x\rho=t\rho\). On variable \(x\), \((\delta\rho)(x)=t\rho=x\rho\);
on every other variable \(v\), \((\delta\rho)(v)=v\rho\).
Thus \(\delta\rho=\rho\) globally. If \(\rho\) also solves \(F\),
the term composition law gives \((F\delta)\rho=F\rho\), so
\(\rho\) solves \(F\delta\) and is a displayed factor \(\delta\rho\).

Conversely, suppose \(\eta\) solves \(F\delta\).
Then \(\delta\eta\) solves \(F\) by the composition law. It solves
the selected binding as well: \(x\delta\eta=t\eta\), while
\(t\delta\eta=t\eta\) because \(x\) is absent from \(t\).
This proves both inclusions of the stated solution-set equality.
The residual substitution may still have an arbitrary image for \(x\);
its value is immaterial to the composition on the forced binding.

For Occurs, select a proper occurrence position of \(x\) in \(t\).
Under any simultaneous substitution \(\eta\), the finite tree
\(x\eta\) occurs at that same proper position in \(t\eta\).
Every constructor above the position remains a constructor. Therefore
\(\operatorname{nodes}(t\eta)>\operatorname{nodes}(x\eta)\),
contradicting literal equality. This does not
assume that \(\eta\) itself is idempotent or acyclic as a finite display.
The case \(t=x\) is excluded precisely because its position is not proper.
\end{proof}
\end{annexitem}

\begin{annexitem}{annex:p12.ue.sheet-invariant}{thm:p12-ue-sheet-invariant}{All six literal transformation cases}
\begin{proof}
Delete removes an equation satisfied by every substitution.
Decompose preserves solutions because in a free ranked term algebra
\(f(s_1,\ldots,s_n)\rho=f(t_1,\ldots,t_n)\rho\) holds exactly
when each corresponding argument image is equal. At arity zero it
removes the identical constant equation and introduces none.
Clash cannot be solved: substitution changes arguments but not a
nonvariable root symbol or its arity. Orient reverses literal term
equality, so has exactly the same solutions; its actual applicability
is restricted to a nonvariable on the left and variable on the right.
Occurs has no solutions by Annex~\ref{annex:p12.ue.forcing}.

For Eliminate, write the old multiset as
\(F\uplus\{x\doteq t\}\) and the new one as
\(F\delta\uplus\{x\doteq t\}\), with the binding retained.
If \(\rho\) solves the old state, the elementary calculation
\(\delta\rho=\rho\) shows that it solves \(F\delta\) as well
as the unchanged binding. Conversely, if \(\rho\) solves the new
state, its retained binding still gives \(\delta\rho=\rho\),
so solution of \(F\delta\) gives solution of \(F\).
These are exactly the same substitutions, including their arbitrary
images on variables absent from the input. The proviso that \(x\)
occur in \(F\) is not needed for this semantic equivalence, but is
essential to the termination/applicability argument for this literal sheet.
Induction through any finite prefix now gives the solution invariant.
\end{proof}
\end{annexitem}

\begin{annexitem}{annex:p12.ue.sheet-termination}{thm:p12-ue-sheet-termination}{Permanent solved variables and a rule-by-rule decreasing triple}
\begin{proof}
Let \(z\in\operatorname{iso}(E)\), with its unique equation
\(z\doteq u\). No rule can select that equation: it is nontrivial,
its left side is a variable, its right side contains no \(z\), and
\(z\) occurs nowhere in the remaining multiset, so Delete, Decompose,
Clash, Orient, Occurs and Eliminate are all disabled on that occurrence.

Rules acting on another equation cannot introduce \(z\).
Delete removes data, Decompose uses only existing subterms, and Orient
only swaps sides. At another elimination \(x\doteq t\),
\(x\ne z\) and \(t\) contains no \(z\), because \(z\) is absent
from that entire selected equation. Substitution may change \(u\)
but introduces no \(z\) there, while its unique left side remains
the variable \(z\). Thus \(z\) remains isolated. This proves
permanence, including updates to earlier solved right sides.
Every transition also introduces no previously absent variable:
Eliminate inserts only variables already in its selected right side;
the other successful rules only delete, select subterms or swap.

For a new Eliminate on \(x\doteq t\), the occurrence proviso says
that \(x\) occurs in the other multiset. Hence \(x\) was not isolated.
After substituting, \(x\) is absent from every other equation and
from \(t\); its retained binding makes it isolated. All old isolated
variables remain isolated, and the variable set does not increase.
Therefore \(\operatorname{uv}\) decreases by at least one. The
possibly large growth of term trees in the second coordinate is irrelevant
in this case; no false total-size monotonicity is assumed.

For Delete, \(\operatorname{uv}\) cannot increase by permanence and
the nonincrease of the variable set. If it stays equal, deleting
\(t\doteq t\) decreases \(\operatorname{sz}\) by
\(2\operatorname{nodes}(t)>0\).
For Decompose, the same first-coordinate argument holds. If it stays
equal, the sum of the child-equation sizes is the sum of the two
argument lists, exactly two less than the selected equation's size,
because the two outer function nodes have disappeared. This includes
zero arity and repeated child equations, counted as a multiset.

For Orient, variables and total size are unchanged. Old isolated
bindings remain isolated; the rule may make a new one isolated, in
which case \(\operatorname{uv}\) decreases. Otherwise its sole
changed equation was counted by \(\operatorname{bad}\) and its
variable-left replacement is not, so that third coordinate decreases
by one. No other occurrence is altered. These are all nonfailure rules.
The lexicographic order on triples of natural numbers is well founded,
so every sequence of actual enabled transitions is finite, regardless
of which equation, variable or applicable rule is chosen.
\end{proof}
\end{annexitem}

\begin{annexitem}{annex:p12.ue.extraction}{thm:p12-ue-extraction}{Every terminal shape and the actual simultaneous answer}
\begin{proof}
Inspect a nonfailure irreducible equation. Equal sides enable Delete.
Two nonvariable sides either have equal root and arity, enabling
Decompose, or clash. A nonvariable left side and variable right side
enable Orient. Thus every remaining equation is \(x\doteq t\)
with a variable on the left and \(x\ne t\). Occurs excludes a
right-side occurrence of \(x\); Eliminate excludes an occurrence of
\(x\) anywhere in the remaining multiset. Consequently the left
variables are distinct, occur exactly once each, and are absent from
all right sides. This also rules out duplicate bindings, right-side
dependencies on a different left variable, and variable cycles.

Let their set be \(X\). The extracted simultaneous substitution
\(\mu\) has support exactly \(X\) and fixes every variable in each
right side. It therefore solves every \(x_i\doteq t_i\), and
\(\mu\mu(x_i)=t_i\mu=t_i=\mu(x_i)\). On variables outside
\(X\), both applications are the identity. This proves idempotence.
No transition introduced a new variable, so \(X\) and every right-side
variable are input variables; the output is identity elsewhere.

By the invariant, \(\mu\) solves the input. If \(\rho\) is any
input solution, it solves the terminal state too. On \(x_i\in X\),
\((\mu\rho)(x_i)=t_i\rho=x_i\rho\); outside \(X\),
\((\mu\rho)(v)=v\rho\). Thus \(\rho=\mu\rho\) on the entire
alphabet, a valid finite-support factorization proving most generality.
Conversely every further instance \(\mu\eta\) remains a solution by
the term composition law. For the empty terminal state, \(X\) is
empty and the same argument gives \(\mu=\mathrm{id}\).

Termination guarantees that every complete run ends either in failure
or in such an irreducible state. Failure has an empty solution set by
the invariant. If the input has an empty solution set, an irreducible
nonfailure state is impossible, since its extracted substitution would
be a solution. Conversely a unifiable input cannot reach a failure
state. This proves both directions for every legal elimination order.
The finite scan described in the body can find an enabled rule or
certify irreducibility; refusing to scan past one solved equation would
not be that complete procedure.
\end{proof}
\end{annexitem}

\begin{annexitem}{annex:p12.ue.accumulator}{thm:p12-ue-accumulator}{The removing procedure, global invariant, range condition and termination}
\begin{proof}
Initially \(E=E_0\), \(\theta=\mathrm{id}\), so the solution-set
invariant holds. Delete, Decompose and Orient preserve the residual
solution set exactly, by the same free-term arguments as before, and
leave \(\theta\) unchanged. Sets may merge identical resulting equations;
duplicates impose no additional solution condition. At an elimination
write the old residual set as \(F\cup\{x\doteq t\}\).
Annex~\ref{annex:p12.ue.forcing} gives
\[
 \{\theta\rho:\rho\in\operatorname{Sol}(F\cup\{x\doteq t\})\}
 =\{\theta(\delta\eta):\eta\in\operatorname{Sol}(F\delta)\}
 =\{(\theta\delta)\eta:\eta\in\operatorname{Sol}(F\delta)\}.
\]
These are global equalities of finite-support substitutions. They are
exactly the new residual and the source's new accumulator. Reversing
that product would not establish this invariant. Clash and Occurs have
no residual solution and therefore no original solution under the invariant.

We now prove the structural accumulator invariant, not merely claim
normalization of an unspecified substitution. Let \(V_0\) be the input
variables and let \(D\) be the variables eliminated so far. Inductively,
\(\operatorname{supp}(\theta)=D\), every right-side variable in
\(\theta(v)\), for \(v\in D\), lies in \(V_0\setminus D\),
and \(\operatorname{Var}(E)\subseteq V_0\setminus D\).
Initially all statements are immediate. Delete, decomposition and
orientation introduce no variables and preserve them. For elimination,
\(x\) belongs to the residual variable set, so \(x\notin D\).
The selected term's variables also lie outside \(D\), and the occurs
test excludes \(x\) from that term. Substituting in the other equations
removes every occurrence of \(x\) and introduces only variables in
\(V_0\setminus(D\cup\{x\})\).

For an old \(v\in D\), the update sends its image to
\((v\theta)\delta\), which contains no variable in
\(D\cup\{x\}\). It cannot become literally \(v\), since that
variable is absent from the updated image. For the new variable \(x\),
the image is \(t\), which differs from \(x\) and has no eliminated
variable. Every variable outside \(D\cup\{x\}\) is still fixed,
because the old accumulator fixes it and \(\delta\) fixes it.
Thus the support becomes exactly \(D\cup\{x\}\), proving the
induction and recording each possible image change. Applying \(\theta\)
again leaves its images unchanged, so \(\theta\theta=\theta\).
In particular, no eliminated variable can be selected again, and no
unused variable acquires a nonidentity image.

For termination use the triple consisting of the number of residual
variables, total sizes of residual term trees, and bad orientations,
now counted over a set. At elimination the selected variable disappears
from the entire residual set, so the first coordinate strictly decreases.
Other transitions do not increase that coordinate. Delete strictly
decreases size if the first coordinate stays equal. Decomposition
removes the two outer nodes; set coalescence can decrease size further,
so its second coordinate strictly decreases when the first is equal.
Orientation preserves size and variables and decreases the last
coordinate by one. These are all successful steps and the lexicographic
triple is well founded. Accumulator growth does not enter this residual
termination measure. A nonempty residual set always offers a rule:
the same outer-term classification used for the sheet ends in a
variable-left equation whose nontrivial binding must be removed here,
even when its variable occurs in no other equation.

At success the residual set is empty. Taking \(\eta=\mathrm{id}\)
in the invariant proves that the returned \(\theta\) is an original
solution. Every other solution is \(\theta\eta\), so it is most general.
Since \(\theta\) is idempotent, such a solution \(\rho=\theta\eta\)
also satisfies \(\theta\rho=\theta(\theta\eta)=\rho\).
An unsolvable input cannot end successfully; a solvable input cannot
fail; termination supplies the remaining dichotomy. An empty input
returns the initial identity substitution without a transition.

Finally, an atom's predicate symbol, or the logical identity atom's
outer sign, cannot be changed by a term substitution. Different outer
symbols/arities therefore cannot be unified. With matching symbols,
literal atom equality is exactly equality of the corresponding argument
terms after the one simultaneous substitution. The preliminary pairing
therefore preserves the solution set; repeated pairs can be merged.
For an identity atom it pairs the two arguments in their written order,
not modulo semantic symmetry. Nullary matching predicates have no
argument equations. This completes the inherited atom-procedure cases.
\end{proof}
\end{annexitem}

\begin{annexitem}{annex:p12.ue.renaming}{thm:p12-ue-renaming}{Constructing the finite variable permutation}
\begin{proof}
Most generality gives finite-support substitutions \(\eta,\zeta\)
with \(\nu=\mu\eta\) and \(\mu=\nu\zeta\), hence
\(\mu=\mu\eta\zeta\). Let the range-variable set of \(\mu\)
be the union of variables in \(v\mu\) over all variables \(v\),
and define the corresponding set for \(\nu\).
If \(w\) occurs in an image \(v\mu\), equality of that finite
term with \((v\mu)\eta\zeta\) makes the term inserted at that
same occurrence position equal to \(w\). Thus \(w\eta\zeta=w\).
The image \(w\eta\) must be a variable: a constant or function
root could not become a variable under substitution. Write it as \(q\).
Then \(q\zeta=w\), so different range variables have different
\(\eta\)-images. Also \(\nu=\mu\eta\) shows that every variable
in the range of \(\nu\) is such an image, and every such image occurs
in that range. Hence \(\eta\) restricts to a bijection between the
two range-variable sets.

Choose a finite variable set \(K\) containing the supports of
\(\mu,\nu,\eta,\zeta\) and all variables in their displayed
nonidentity images. Every variable outside \(K\) is fixed by these
substitutions, lies in both range-variable sets, and is fixed by the
restricted bijection. The bijection maps the first range set's part
inside \(K\) onto the second one's part inside \(K\): images of
members of \(K\) lie in \(K\), and outside members stay outside.
These finite parts have equal cardinality, so their complements in
\(K\) do as well. Pair the two complements using any finite bijection,
for example their induced variable-alphabet order. Together with the
range-set bijection inside \(K\) and the identity outside \(K\),
this defines a finite-support permutation \(\pi\) of all variables.
It agrees with \(\eta\) on every variable in every \(\mu\)-image,
so the term induction gives \(\mu\pi=\mu\eta=\nu\).

Conversely, \(\mu\pi\) remains a solution because term renaming
preserves each solved equation. If \(\rho=\mu\delta\) is any
solution, then \(\rho=(\mu\pi)(\pi^{-1}\delta)\), a finite-support
factorization, so \(\mu\pi\) is most general. This proves the
answer comparison, without asserting uniqueness of either \(\eta\)
or \(\delta\).
\end{proof}
\end{annexitem}

\begin{annexitem}{annex:p12.ue.comparison}{thm:p12-ue-comparison}{Duplicate constraints and answer equivalence do not identify the procedures}
\begin{proof}
Passing from a finite equation multiset to its support set preserves
exactly the solutions, since repeated requirements assert the same
literal equality. The two proved answer theorems therefore classify
the same input as solvable or not. In the solvable case both outputs
are finite-support most general solutions in the same global instance
preorder. Apply Annex~\ref{annex:p12.ue.renaming} to obtain the finite
variable permutation relating them. The argument also compares two
different legal runs of either procedure. It uses their solution and
most-generality guarantees, not identity of intermediate states or
an invented correspondence between individual transition traces.
\end{proof}
\end{annexitem}
